% ======================================================================
\section{Task and metric}\label{sec:task}

\subsection{Terminology}\label{sec:terms}
\begin{description}[itemsep=1pt]
  \item[Final system] the system described in this report (release v1.2.0 of the code).
  \item[Direct audio-to-image] the baseline that draws every detected sound without checking whether its source is on
    screen (Section~\ref{sec:baselines}).
  \item[Show nothing] a system that never shows a picture.
  \item[Gate (on-screen check)] the step that decides whether a sound's source is visible.
  \item[Listener] an audio-language model asked about a sound (``is a siren present?'', ``list what you hear'').
  \item[Family] two labels are in the same family if they are equal or one lies under the other in the AudioSet tree
    (\emph{Bark} and \emph{Dog}), and neither is a vague label near the top of the tree such as \emph{Animal}.
  \item[Span] the time interval of one detected sound.
  \item[Stretch] a piece of a span, about 5\,s long; the on-screen check judges each stretch separately.
  \item[Bootstrap] resampling the clips many times to see how much a difference could change by chance (for example,
        draw 88 clips with replacement 10{,}000 times and recompute the cost difference each time).
  \item[Holm correction] a stricter $p$-value bar used when several comparisons are made at once.
\end{description}


\subsection{Needed sounds}\label{sec:needed}
The unit of evaluation is one sound in one clip. For every sound the annotator records its label, its start and end
(about 0.5-s precision) and three judgements:
\begin{description}
  \item[visible] the thing making the sound is \emph{seen making it now}. A barking dog on screen is visible; a baby
        held quietly is not visibly crying; an alarm box on a wall is not visibly ringing.
  \item[obvious] with the sound off, a viewer would still know it is happening (waves when the sea fills the frame).
  \item[importance] rated as if the screen were black: 1 = steady background noise; 2 = an event one can say in a
        sentence; 3 = danger or a key moment (siren, alarm, gunshot, glass, baby crying, scream).
\end{description}
A sound is \textbf{needed} if it is neither visible nor obvious and has importance 2 or 3. \emph{Example}: a police
siren is heard while the camera shows a crowd looking up -- needed. When the police car drives into view with its
lights on, the siren becomes visible -- not needed. Every heard sound is labelled, because a picture over an unlabelled
moment counts as wrong.

\subsection{Scoring a picture}\label{sec:matching}
Pictures are scored as they are shown on screen. A picture \textbf{hits} a needed sound if it is of the same family
and \emph{starts} between 0.5\,s before and 1\,s after the sound starts; each sound takes at most one picture. Every
other picture is \textbf{wrong}, of one of three kinds:
\begin{description}
  \item[on screen] it shows a sound whose source is visible (the viewer could see it);
  \item[other sound] it shows the wrong sound -- something else is heard there (a gunshot picture over footsteps), or
        it is a late second picture of a long sound already shown;
  \item[nothing] there is no labelled sound there at all.
\end{description}
``Wrong'' is therefore strict: it includes redundant pictures, not only false ones.

\subsection{Viewer cost}\label{sec:metric}
The main score is the cost per clip,
\begin{equation}
  C \;=\; \frac{4\cdot\#\text{missed needed sounds} \;+\; 2\cdot\#\text{wrong pictures}}{\#\text{clips}},
  \label{eq:cost}
\end{equation}
lower is better. A miss costs twice a wrong picture, so one right picture pays for two wrong ones, and a picture is
worth showing when it is right at least one time in three. These weights are an assumption, not a measurement: they
were chosen early in the project for an earlier comparison, and became the main score after an earlier test found no
significant difference in F1 between systems (F1 counts a miss and a wrong picture alike). Section~\ref{sec:sensitivity} therefore also reports other weights, including a version in which an on-screen
picture costs only 1 or 0 (it is less harmful than a picture of nothing):
\begin{equation}
  C_w \;=\; \frac{4\cdot\#\text{missed} + w\cdot\#\text{on screen} + 2\cdot\#\text{other sound} + 2\cdot\#\text{nothing}}{\#\text{clips}},
  \qquad w\in\{0,1,2\}.
  \label{eq:costw}
\end{equation}
\emph{Why not F1?} F1 counts a miss and a wrong picture the same, while the gate's value lies almost entirely in the
wrong pictures it removes. F1 is reported alongside the cost (Table~\ref{tab:headline}); the cost is primary.

\subsection{How decisions were made}\label{sec:passrules}
Comparisons are paired by clip and tested with a paired clip bootstrap. A change to the system was accepted only if
it passed a rule written down \emph{before} its run, on the development set (the rule and its pass bar were
committed to the repository first; this is called pre-registration). It passed if either
\begin{enumerate}[label=(\alph*)]
  \item it kept or raised the hits, added at most two wrong pictures per extra hit, and lowered the cost; or
  \item it lost at most three hits and removed at least three wrong pictures per hit lost.
\end{enumerate}
Rules about the audio alone could first be screened on the screening set (Section~\ref{sec:benchmark}). The test set
was scored only after a change was accepted, and only reported; no decision in the final system used it. The 95\,\% intervals of the main comparisons are 0.5 to 1.2 per clip wide on both sets, so small single changes could only be judged on the development rules.
